\section*{Chapitre \ref{ch:b3:rovibrational-spectroscopy} --- Spectroscopies de rotation et de vibration}

\begin{solution}{exo:b3:rovibrational-spectroscopy:1}
$I = h/(8\pi^2cB) = 6.62607 \times 10^{-34}/(8\pi^2 \times 2.99792 \times 10^{10} \times
1.9225) = \qty{1.4561e-46}{kg.m^2}$. $\mu = \qty{6.85621}{u} = \qty{1.13850e-26}{kg}$, donc
$r_0 = \sqrt{I/\mu} = \qty{113.09}{pm}$.
\end{solution}

\begin{solution}{exo:b3:rovibrational-spectroscopy:2}
$2B_0(J+1)$~: 20.88, 41.76, 62.64, \qty{83.52}{cm^{-1}}. $J_{\max} = \sqrt{kT/2hcB} -
\frac12 = \sqrt{208.5/20.88} - 0.5 = 2.7$~: le niveau $J = 3$ est le plus
peuplé, et la raie la plus intense est $4 \leftarrow 3$, à
\qty{83.5}{cm^{-1}}.
\end{solution}

\begin{solution}{exo:b3:rovibrational-spectroscopy:3}
L'absorption micro-ondes exige un dipôle permanent~: \ce{HCl} et \ce{H2O}
seulement. Le Raman de rotation exige une polarisabilité anisotrope~:
\ce{H2}, \ce{HCl}, \ce{N2},
\ce{CO2}, \ce{H2O}~; pas les \omterm{def:b3:rovibrational-spectroscopy:rotor-types}{toupies sphériques} \ce{CH4} et \ce{SF6}.
\end{solution}

\begin{solution}{exo:b3:rovibrational-spectroscopy:4}
$N_1/N_0 = \eu^{-hc\tilde\nu/kT}$~: à \qty{300}{K}, $\eu^{-13.84} = \num{9.8e-7}$~; à
\qty{1000}{K}, $\eu^{-4.15} = 0.016$. La \omterm{def:b3:rovibrational-spectroscopy:anharmonic}{bande chaude} ne devient sensible
(un pour cent de la fondamentale) que vers \qty{1000}{K}.
\end{solution}

\begin{solution}{exo:b3:rovibrational-spectroscopy:5}
$\tilde\omega_e - 2\tilde\omega_ex_e = 2885.3$ et $2\tilde\omega_e - 6\tilde\omega_ex_e =
5665.0$~: en retranchant le double de la première de la seconde, $-2\tilde\omega_ex_e = -105.6$,
donc $\tilde\omega_ex_e = \qty{52.8}{cm^{-1}}$ et $\tilde\omega_e = \qty{2990.9}{cm^{-1}}$.
\end{solution}

\begin{solution}{exo:b3:rovibrational-spectroscopy:6}
$D_e = 2990.95^2/(4 \times 52.82) = \qty{42341}{cm^{-1}}$~; $G(0) = 1495.47 - 13.20 =
\qty{1482.3}{cm^{-1}}$~; $D_0 = \qty{40859}{cm^{-1}} = \qty{488.8}{kJ/mol}$,
contre la vraie valeur
\qty{35760}{cm^{-1}} $= \qty{427.8}{kJ/mol}$~: le modèle de Morse
surestime de 14\,\%.
\end{solution}

\begin{solution}{exo:b3:rovibrational-spectroscopy:7}
$R(0) - P(2) = 6B_0 = 62.64$~: $B_0 = \qty{10.440}{cm^{-1}}$. $R(1) - P(1) = 6B_1 = 60.80$~:
$B_1 = \qty{10.133}{cm^{-1}}$. $\alpha_e = \qty{0.307}{cm^{-1}}$. $R(0) = \tilde\nu_0 + 2B_1$ donne
$\tilde\nu_0 = \qty{2885.31}{cm^{-1}}$ (et $P(1) = \tilde\nu_0 - 2B_0$
concorde).
\end{solution}

\begin{solution}{exo:b3:rovibrational-spectroscopy:8}
$\mu_H = \qty{0.97959}{u}$, $\mu_D = \qty{1.90441}{u}$, $\mu_H/\mu_D = 0.51438$. $B_0(\ce{DCl})
= 10.440 \times 0.51438 = \qty{5.370}{cm^{-1}}$. $\tilde\omega_e = 2990.95 \times \sqrt{0.51438}
= 2145.1$, $\tilde\omega_ex_e = 52.82 \times 0.51438 = 27.17$: $\tilde\nu_0 = 2145.1 - 54.3 =
\qty{2090.8}{cm^{-1}}$.
\end{solution}

\begin{solution}{exo:b3:rovibrational-spectroscopy:9}
Déplacements $4B_0(J + \frac32)$~: 11.94, 19.90, \qty{27.85}{cm^{-1}}
(depuis $J = 0$, 1, 2). Les
deux noyaux \ce{^{14}N} sont des bosons identiques de spin 1~: les
niveaux de $J$ pair ont six états de spin nucléaire, ceux de $J$ impair
trois, si bien que les raies issues de $J$ pair sont deux fois plus
intenses que leurs voisines.
\end{solution}

\begin{solution}{exo:b3:rovibrational-spectroscopy:10}
Un \omterm{def:b3:quantum-model-systems:rotor}{rotateur rigide} placerait les raies à $\nu_1$, $2\nu_1$, $3\nu_1$~: $230542.40$ et
\qty{345813.61}{MHz}, au-dessus des valeurs mesurées de 4.4 et
\qty{17.6}{MHz}. Avec
$\nu_{J+1\leftarrow J} = 2B(J+1) - 4D(J+1)^3$~: $\nu_1 = 2B - 4D$, $\nu_3 = 6B - 108D$, donc
$3\nu_1 - \nu_3 = 96D$~: $D = \qty{0.1835}{MHz}$ et $B = (\nu_1 + 4D)/2 =
\qty{57635.97}{MHz}$. Vérification~: $\nu_2 = 4B - 32D = \qty{230538.0}{MHz}$.
\end{solution}

\begin{solution}{exo:b3:rovibrational-spectroscopy:11}
$\Delta G(v + \frac12) = \tilde\omega_e - 2\tilde\omega_ex_e(v+1)$ s'annule vers $v =
\tilde\omega_e/2\tilde\omega_ex_e - 1$~; la somme d'une progression
arithmétique de $n$ termes
allant de $\tilde\omega_e - 2\tilde\omega_ex_e$ à presque 0 vaut environ $n\tilde\omega_e/2 \approx
\tilde\omega_e^2/4\tilde\omega_ex_e - \tilde\omega_e/2 \approx D_e - G(0)$. Pour \ce{HCl}, les 28 écarts
($v = 0$ à 27) ont pour somme \qty{40857}{cm^{-1}}, contre
$D_e - G(0) = \qty{40859}{cm^{-1}}$.
\end{solution}

\begin{solution}{exo:b3:rovibrational-spectroscopy:12}
Avec les seuls $J$ pairs, les transitions $J \to J + 2$ partent de
$J = 0, 2, 4, \dots$~: déplacements
$6B, 14B, 22B, \dots$, espacement $8B$. \ce{CO2} possède un centre de
symétrie et pas de dipôle permanent~: pas de spectre micro-ondes.
\end{solution}

\begin{solution}{pb:b3:rovibrational-spectroscopy:1}
\textbf{1.} $B_0 = \nu/2 = \qty{57.6356}{GHz}$~; $B_0/c = \qty{1.92252}{cm^{-1}}$.
\textbf{2.} $\mu = 12 \times 15.99491/27.99491 = \qty{6.85621}{u} = \qty{1.13850e-26}{kg}$.
\textbf{3.} $I = h/8\pi^2B_0 = \qty{1.4560e-46}{kg.m^2}$.
\textbf{4.} $r_0 = \sqrt{I/\mu} = \qty{113.09}{pm}$.
\textbf{5.} $r_0$ dépasse $r_e = \qty{112.83}{pm}$ de \qty{0.26}{pm}~:
$B_0$ moyenne $1/r^2$
sur la vibration de point zéro dans un puits anharmonique, qui passe
plus de temps aux grandes distances.
\textbf{6.} $2\nu_{1\leftarrow0} = \qty{230.5424}{GHz}$~; la raie mesurée
est plus basse de \qty{4.4}{MHz}~:
distorsion centrifuge.
\textbf{7.} $0$, 3.845 et \qty{11.535}{cm^{-1}}, soit $E/k = 0$, 5.53 et \qty{16.60}{K}.
\textbf{8.} $N_2/N_1 = \frac53\exp[-(16.60 - 5.53)\,\mathrm K/T]$.
\textbf{9.} $\ln(\frac53/0.85) = 0.673 = 11.06/T$~: $T = \qty{16}{K}$.
\textbf{10.} $J_{\max} = \sqrt{16.4/(2 \times 1.43878 \times 1.9225)} - 0.5 = 1.2$~: $J = 1$.
\textbf{11.} Un quantum de vibration correspond à $hc\tilde\nu/k \approx \qty{3080}{K}$~: à
\qty{16}{K}, aucune molécule n'est excitée en vibration, si bien que le
nuage n'émet aucune lumière de vibration.
\textbf{12.} Elle ne demande que \qty{5.5}{K} d'excitation, est excitée
même dans le gaz le plus froid, et CO est abondant et possède un moment
dipolaire.
\textbf{13.} $\mu = 13.00335 \times 15.99491/28.99826 = \qty{7.17241}{u}$.
\textbf{14.} $B \propto 1/\mu$~: $115.2712 \times 6.85621/7.17241 = \qty{110.1893}{GHz}$.
\textbf{15.} La prévision est inférieure de \qty{12}{MHz} (0.011\,\%) à
la raie mesurée~: $B_0 = B_e - \frac12\alpha_e$, et le terme de
vibration--rotation $\alpha_e$ varie selon une autre puissance de
$\mu$~; la structure d'équilibre est la même, la moyenne de point zéro ne
l'est pas.
\textbf{16.} $98.9/1.1 = 90$.
\textbf{17.} La raie de \ce{^{12}C^{16}O} est saturée (optiquement
épaisse)~: elle ne croît plus avec la quantité de gaz.
\textbf{18.} La raie du rare \ce{^{13}C^{16}O} reste proportionnelle à
la quantité de gaz, la raie de
\ce{^{12}C^{16}O} donne la température~: ensemble, elles mesurent les
deux.
\textbf{19.} $\tilde\nu_0 = 2169.81 - 2 \times 13.288 = \qty{2143.24}{cm^{-1}}$, \qty{4.666}{\micro m}.
\textbf{20.} $2\tilde\omega_e - 6\tilde\omega_ex_e = \qty{4259.89}{cm^{-1}}$.
\textbf{21.} $G(0) = \frac12\tilde\omega_e - \frac14\tilde\omega_ex_e = \qty{1081.58}{cm^{-1}}$.
\textbf{22.} $R(0) - P(1) = 2(B_0 + B_1) \approx 4B \approx \qty{7.7}{cm^{-1}}$.
\textbf{23.} \ce{^{13}C^{16}O}~: \omterm{def:b3:quantum-model-systems:oscillator}{masse réduite} plus grande, fréquence plus
basse.
\textbf{24.} \textbf{$r_0(\ce{CO}) = \qty{113.1}{pm}$}, tiré d'une seule
fréquence radio.
\end{solution}
